% ================================================================
% 📚 GUÍA DE USO DE COMANDOS PERSONALIZADOS
% ================================================================
% Este archivo muestra cómo usar todos los comandos disponibles
% en template_pid.cls
% ================================================================

\documentclass{template_pid}

\begin{document}

\chapter{Guía de Comandos Personalizados}

\section{Formato de Texto Básico}

\subsection{Cursiva, Negrita y Subrayado}

Puedes usar \italica{texto en cursiva} para énfasis suave.

Usa \negrita{texto en negrita} para dar más peso.

El \subrayado{texto subrayado} también está disponible.

Para combinar estilos: \negrita{\italica{negrita e itálica juntas}}.

\subsection{Términos y Conceptos Importantes}

Define un \termino{término técnico} importante (negrita + cursiva).

Destaca un \concepto{concepto clave} en color azul.

El comando \importante{¡IMPORTANTE!} resalta en rojo y negrita.

\subsection{Código y Variables}

Para código inline usa \codigo{printf("Hola mundo")}.

Para variables o parámetros: \variable{nombreVariable}.

Las \sigla{UCI}, \sigla{HTML}, \sigla{CSS} se escriben en versalitas.

\subsection{Colores Personalizados}

Puedes usar \colortext{green}{texto en verde} o \colortext{purple}{texto en morado}.

% ================================================================
\section{Notas y Observaciones}

\nota{Nota al margen}
Este texto tiene una nota al margen. Las notas aparecen en el margen derecho.

\notaimportante{Esta es una nota importante que aparece destacada en el texto principal en color naranja.}

\observacion{Esta es una observación relevante que se muestra en color teal y cursiva.}

% ================================================================
\section{Entornos Personalizados}

\subsection{Ejemplos}

\begin{ejemplo}
Este es un ejemplo de cómo usar un comando o concepto.
Puede contener múltiples líneas y explicaciones detalladas.
\end{ejemplo}

\subsection{Definiciones}

\begin{definicion}{Base de Datos}
Una base de datos es un conjunto organizado de datos estructurados,
normalmente almacenados electrónicamente en un sistema informático.
\end{definicion}

\subsection{Bloques de Código}

\begin{codigobloque}
def suma(a, b):
    return a + b

resultado = suma(5, 3)
print(resultado)  # Output: 8
\end{codigobloque}

\subsection{Advertencias}

\begin{advertencia}
Este proceso puede consumir mucha memoria. Asegúrate de tener
al menos 8GB de RAM disponibles antes de ejecutar.
\end{advertencia}

% ================================================================
\section{Listas y Enumeraciones}

\subsection{Lista con Números}

\begin{enumeracion}
    \item Primer elemento de la lista
    \item Segundo elemento de la lista
    \item Tercer elemento de la lista
    \begin{enumeracion}
        \item Subelemento 1
        \item Subelemento 2
    \end{enumeracion}
    \item Cuarto elemento
\end{enumeracion}

\subsection{Lista con Viñetas}

\begin{listapuntos}
    \item Primer punto
    \item Segundo punto
    \item Tercer punto
    \begin{listapuntos}
        \item Subpunto A
        \item Subpunto B
    \end{listapuntos}
\end{listapuntos}

\subsection{Lista Alfabética}

\begin{listaalfabetica}
    \item Primera opción (a)
    \item Segunda opción (b)
    \item Tercera opción (c)
    \item Cuarta opción (d)
\end{listaalfabetica}

\subsection{Lista con Números Romanos}

\begin{listaromana}
    \item Fase I del proyecto
    \item Fase II del proyecto
    \item Fase III del proyecto
    \item Fase IV del proyecto
\end{listaromana}

% ================================================================
\section{Referencias Cruzadas}

Supongamos que tienes una sección etiquetada como \codigo{\textbackslash label\{sec:introduccion\}}

Puedes referenciarla con: Ver \refsec{sec:introduccion} para más detalles.

Para figuras: Como se muestra en \reffig{fig:arquitectura}...

Para tablas: Los datos se presentan en \reftab{tab:resultados}...

Para epígrafes: Consultar \refepi{epi:metodologia}...

% ================================================================
\section{Ejemplos Combinados}

\subsection{Caso de Uso Práctico}

En el development de \termino{sistemas web}, es fundamental comprender
el \concepto{patrón MVC} (Modelo-Vista-Controlador).

\begin{definicion}{MVC}
\sigla{MVC} es un patrón de arquitectura de software que separa
una aplicación en tres componentes principales:
\begin{listapuntos}
    \item \negrita{Modelo}: maneja los datos y la lógica de negocio
    \item \negrita{Vista}: presenta la información al usuario
    \item \negrita{Controlador}: gestiona la interacción del usuario
\end{listapuntos}
\end{definicion}

\begin{ejemplo}
En Django (Python), puedes definir un modelo así:

\begin{codigobloque}
from django.db import models

class Usuario(models.Model):
    nombre = models.CharField(max_length=100)
    email = models.EmailField(unique=True)
    fecha_registro = models.DateTimeField(auto_now_add=True)
\end{codigobloque}
\end{ejemplo}

\notaimportante{Es importante validar siempre los datos de entrada
antes de guardarlos en la base de datos.}

\subsection{Proceso de Implementación}

Para implementar un sistema \sigla{CRUD} básico:

\begin{enumeracion}
    \item Diseñar el modelo de datos
    \item Crear las vistas necesarias
    \item Implementar los controladores
    \item Realizar pruebas unitarias
    \item Desplegar en producción
\end{enumeracion}

\begin{advertencia}
Nunca subas las credenciales de base de datos al repositorio.
Usa variables de entorno para información sensible.
\end{advertencia}

% ================================================================
\section{Resumen de Comandos Disponibles}

\subsection{Formato de Texto}
\begin{listapuntos}
    \item \codigo{\textbackslash italica\{texto\}} - Texto en cursiva
    \item \codigo{\textbackslash negrita\{texto\}} - Texto en negrita
    \item \codigo{\textbackslash subrayado\{texto\}} - Texto subrayado
    \item \codigo{\textbackslash termino\{texto\}} - Término importante (negrita + cursiva)
    \item \codigo{\textbackslash concepto\{texto\}} - Concepto destacado (azul + negrita)
    \item \codigo{\textbackslash importante\{texto\}} - Texto muy importante (rojo + negrita)
    \item \codigo{\textbackslash sigla\{ABC\}} - Siglas en versalitas
\end{listapuntos}

\subsection{Código}
\begin{listapuntos}
    \item \codigo{\textbackslash codigo\{codigo\}} - Código inline
    \item \codigo{\textbackslash variable\{var\}} - Variable o parámetro
\end{listapuntos}

\subsection{Notas}
\begin{listapuntos}
    \item \codigo{\textbackslash nota\{texto\}} - Nota al margen
    \item \codigo{\textbackslash notaimportante\{texto\}} - Nota destacada (naranja)
    \item \codigo{\textbackslash observacion\{texto\}} - Observación (teal + cursiva)
\end{listapuntos}

\subsection{Entornos}
\begin{listapuntos}
    \item \codigo{ejemplo} - Para mostrar ejemplos
    \item \codigo{definicion\{nombre\}} - Para definiciones formales
    \item \codigo{codigobloque} - Para bloques de código
    \item \codigo{advertencia} - Para advertencias importantes
\end{listapuntos}

\subsection{Listas}
\begin{listapuntos}
    \item \codigo{enumeracion} - Lista numerada (1, 2, 3...)
    \item \codigo{listapuntos} - Lista con viñetas (•)
    \item \codigo{listaalfabetica} - Lista con letras (a, b, c...)
    \item \codigo{listaromana} - Lista romana (I, II, III...)
\end{listapuntos}

\subsection{Referencias}
\begin{listapuntos}
    \item \codigo{\textbackslash refsec\{label\}} - Referencia a sección
    \item \codigo{\textbackslash reffig\{label\}} - Referencia a figura
    \item \codigo{\textbackslash reftab\{label\}} - Referencia a tabla
    \item \codigo{\textbackslash refepi\{label\}} - Referencia a epígrafe
\end{listapuntos}

\end{document}
